\begin{textsample}{POS Dim 2 – 2017\_c6 – Score -1.00 – 2017\_c6/enade\_2017\_eng\_contr\_4e26513cc7.txt}  \label{ex:f2_pos_271}
CAPÍTULO 6 ANÁLISE TÉCNICA DA PROVA Este capítulo tem por objetivo apresentar o desempenho dos estudantes concluintes de Engenharia de Controle e Automação no Enade/2017.
Para isso, foram calculadas as estatísticas básicas da prova em seu todo ( seção 6.1.1 ), bem como as estatísticas relacionadas aos Componentes de Formação Geral ( seção 6.1.2 ) e de Conhecimento Específico da Área ( seção 6.1.3 ).
Dadas as suas características, foram analisadas, em separado, as questões objetivas ( seção 6.2 ) e as questões discursivas ( seção 6.3 ).
Para as questões objetivas, foram disponibilizados os índices de facilidade e de discriminação ponto-bisserial, também em separado, para os Componentes de Formação Geral ( seção 6.2.1 ) e de Conhecimento Específico ( seção 6.2.2 ).
De cada componente, uma das questões foi escolhida para exemplificar a análise gráfica, relacionando as alternativas escolhidas pelos estudantes ( inclusive o gabarito ) com o número de acertos no componente.
O Anexo I apresenta a íntegra da análise gráfica para todas as questões objetivas.
Para cada uma das questões discursivas, os conteúdos dos tipos mais comuns de respostas dos estudantes são apresentados e comparados com o padrão de respostas esperado ( ver Anexo VIII com o padrão de respostas ).
Tomando como base as duas questões discursivas do Componente de Formação Geral, a seção 6.3.1.6 apresenta comentários sobre a correção das respostas com respeito à Língua Portuguesa.
As tabelas apresentam as seguintes estatísticas das notas¹⁹ : média do desempenho na prova, erro padrão da média, desvio padrão, nota mínima, mediana e nota máxima para cada um de seus componentes.
Tais estatísticas contemplam o total de estudantes concluintes da Área de Engenharia de Controle e Automação inscritos e presentes à prova do Enade/2017, tendo em vista agregações, ou por Grandes Regiões e o país como um todo, ou por Categoria Administrativa, Organização Acadêmica da IES e Modalidade de Ensino.
Em relação aos gráficos de distribuição de notas, o intervalo considerado foi de 10 unidades, aberto à esquerda e fechado à direita, com exceção do primeiro intervalo, [ 0 ; 10 ], fechado em ambos os extremos.
Para os gráficos de distribuição das notas das questões discursivas, foram consideradas mais duas categorias : questão em branco²⁰ e nota zero. ¹⁹ Uma definição dessas estatísticas pode ser encontrada no Glossário. ²⁰ Nesse grupo estão incluídas também as respostas classificadas como nulas ou desconsideradas. 6.1 ESTATÍSTICAS BÁSICAS DA PROVA Esta seção apresenta estatísticas selecionadas e histogramas da nota geral ( 6.1.1 ) e de cada componente : Formação Geral ( 6.1.2 ) e Conhecimento Específico ( 6.1.3 ).
São também apresentadas estatísticas selecionadas de subpopulações caracterizadas por Grande Região, Categoria Administrativa, Organização Acadêmica e Modalidade de Ensino. 6.1.1 Estatísticas Básicas Gerais A Tabela 6.1 apresenta as Estatísticas Básicas da prova, por Grande Região, dos estudantes concluintes de Engenharia de Controle e Automação.
A Média das notas da prova como um todo ( nas seções seguintes serão analisados os Componentes de Formação Geral e de Conhecimento Específico ) foi 42,1, sendo que os alunos da região Sul obtiveram a Média mais alta ( 48,2 ), e os da região Norte obtiveram a Média mais baixa ( 38,8 ).
As demais médias foram : 41,2 na região Sudeste e 42,2 nas regiões Nordeste e Centro-Oeste,.
O Desvio padrão para o Brasil como um todo foi 13,2, sendo o maior Desvio padrão encontrado na região Sul ( 13,7 ) e o menor, na região Norte ( 10,6 ), indicando uma dispersão um pouco menor das notas desta última região.
A região que obteve a maior nota Máxima foi a Sudeste ( 91,0 ), ao passo que a região que atingiu a menor nota Máxima foi a Norte ( 72,3 ).
A Mediana do Brasil como um todo foi 41,5, sendo a maior Mediana obtida na região Sul ( 48,4 ), e a menor obtida na Norte ( 39,1 ).
A nota Mínima foi zero em duas regiões, Nordeste e Sudeste.
Nas demais, a nota Mínima foi 12,1 na região Norte, 13,8 na região Sul e 13,0 na região Centro-Oeste.
Considerando -se as notas segundo Grande Região, observa -se que existe diferença estatisticamente significativa ao nível de 95 \% entre a maior Média, obtida na região Sul ( 48,2 ), e as médias de todas as demais regiões²¹. ²¹ Todas as comparações deste capítulo utilizam os intervalos de 95 \%.
Os erros-padrão da média que possibilitam os testes estão disponíveis nas tabelas.
O comportamento das notas dos estudantes de todo o Brasil pode ser observado no Gráfico 6.1 que apresenta um histograma com a distribuição das mesmas.
Essa é uma distribuição unimodal com a moda no intervalo ( 40 ; 50 ], seguido pelo intervalo ( 30 ; 40 ].
A Tabela 6.2 apresenta informações referentes à Média da nota final desagregadas por Categoria Administrativa, por Organização Acadêmica e por Modalidade de Ensino.
A Média das notas da prova como um todo foi 42,1.
Em relação à Categoria Administrativa, os estudantes das IES Públicas obtiveram Média mais alta ( 47,4 ), e os das IES Privadas obtiveram Média mais baixa ( 38,0 ), que a Média nacional.
Observa -se que existe diferença estatisticamente significativa entre as médias das notas das IES Públicas e Privadas.
A diferença entre as médias das regiões Norte e Sul ( 9,4 ), a maior e a menor Média, é igual à diferença entre IES Públicas e Privadas ( 9,4 ), caracterizando a não existência de uma maior diversidade administrativa ou regional.
O Desvio padrão para as IES Públicas ( 13,1 ) foi praticamente o mesmo do Brasil como um todo ( 13,2 ), indicando uma dispersão semelhante à do Brasil nesta Categoria Administrativa.
No tocante à Organização Acadêmica, as Universidades e os CEFET/IFET obtiveram Média mais alta que a nacional ( 45,5, 44,5 e 42,1, respectivamente ).
A Média das Faculdades ( 36,2 ) e dos Centros Universitários ( 38,5 ) foram menores.
Constata -se que existe diferença estatisticamente significativa ao nível de 95 \% entre as Médias mais baixas, das Faculdades e dos Centros universitários, e a dos outros dois tipos de Organização Acadêmica.
Quanto à Modalidade de Ensino, não houve estudantes desta área na modalidade a distância.
A Média da Educação Presencial foi 42,1, a mesma da prova como um todo. 6.1.2 Estatísticas Básicas no Componente de Formação Geral A Tabela 6.3 apresenta as Estatísticas Básicas em relação ao componente da prova que avalia a Formação Geral dos estudantes concluintes.
Os alunos de todo o Brasil obtiveram desempenho médio de 57,4.
Quanto à variabilidade, o Desvio padrão das notas dos estudantes do Brasil, como um todo, foi 17,9.
A maior Média foi obtida na região Sul ( 61,6 ), e a menor, na região Norte ( 56,4 ).
As demais médias foram : 57,8 na região Nordeste, 56,7 na região Sudeste e 56,9 na A maior nota no Componente de Formação Geral da prova do Enade/2017 foi 98,2, obtida por, pelo menos, um aluno na região Centro-Oeste.
A menor nota Máxima foi obtida na região Norte ( 93,4 ).
A Mediana do Brasil, como um todo, foi 58,8, sendo a menor Mediana encontrada na região Norte ( 56,8 ), e a maior, encontrada na região Sul ( 63,1 ).
A nota Mínima nessa parte foi zero nas regiões Nordeste e Sudeste, 7,5 nas regiões Sul e Centro-Oeste e 15,0 na região Norte.
Considerando -se as notas segundo Grande Região, observa -se que existe diferença estatisticamente significativa entre a maior Média das notas do Componente de Formação Geral, obtida na região Sul ( 61,6 ), e as demais médias regionais.
O Gráfico 6.2 propicia a avaliação do desempenho dos estudantes do Componente de Formação Geral, a partir do histograma da distribuição das notas correspondentes.
A distribuição é unimodal, com moda em ( 50 ; 60 ], seguida de perto pelo intervalo ( 60 ; 70 ]. região Centro-Oeste.
Já o maior Desvio padrão foi obtido nas regiões Sudeste e Centro-Oeste ( 18,1 em ambas ), e o menor, na região Norte ( 15,1 ).
Os demais desvios padrões foram : 17,8 na região Nordeste e 17,1 na região Sul.
A maior nota no Componente de Formação Geral da prova do Enade/2017 foi 98,2, obtida por, pelo menos, um aluno na região Centro-Oeste.
A menor nota Máxima foi obtida na região Norte ( 93,4 ).
A Mediana do Brasil, como um todo, foi 58,8, sendo a menor Mediana encontrada na região Norte ( 56,8 ), e a maior, encontrada na região Sul ( 63,1 ).
A nota Mínima nessa parte foi zero nas regiões Nordeste e Sudeste, 7,5 nas regiões Sul e Centro-Oeste e 15,0 na região Norte.
Considerando -se as notas segundo Grande Região, observa -se que existe diferença estatisticamente significativa entre a maior Média das notas do Componente de Formação Geral, obtida na região Sul ( 61,6 ), e as demais médias regionais.
O Gráfico 6.2 propicia a avaliação do desempenho dos estudantes do Componente de Formação Geral, a partir do histograma da distribuição das notas correspondentes.
A distribuição é unimodal, com moda em ( 50 ; 60 ], seguida de perto pelo intervalo ( 60 ; 70 ].
Na Tabela 6.4 são apresentadas as informações referentes ao desempenho dos concluintes do Componente de Formação Geral, em diferentes agregações : Categoria Administrativa, Organização Acadêmica e Modalidade de Ensino.
Observa -se que existe diferença estatisticamente significativa entre as médias dos tipos de Categoria Administrativa : a maior Média obtida por estudantes de IES Públicas ( 62,5 ) e a menor pelos de IES Privadas ( 53,4 ).
Nota -se que também há diferença estatisticamente significativa entre a menor Média, obtida em Faculdades ( 51,3 ), e as médias dos demais tipos de Organização Acadêmica.
No que se refere à Modalidade de Ensino, como mencionado anteriormente, só houve estudantes de Educação Presencial, cuja Média foi 57,4. 6.1.3 Estatísticas Básicas do Componente de Conhecimento Específico A Tabela 6.5 apresenta as Estatísticas Básicas referentes ao Componente de Conhecimento Específico da área de Engenharia de Controle e Automação.
A Média do desempenho dos alunos do Brasil, como um todo, foi 37,0.
A maior Média foi obtida na região Sul ( 43,7 ), e a menor, na região Norte ( 32,9 ).
As demais médias foram : 37,0 na região Nordeste, 36,0 na região Sudeste e 37,3 na região Centro-Oeste.
Quanto à variabilidade das notas, o Desvio padrão do Brasil, como um todo, foi 14,3, sendo o maior Desvio padrão observado na região Sul ( 15,4 ), e o menor, na região Norte ( 11,6 ).
Os demais desvios foram : 13,1 na região Nordeste e 14,1 nas regiões Sudeste e Centro-Oeste.
A Mediana das notas dos estudantes de todo o Brasil foi 35,7.
A maior Mediana ocorreu na região Sul ( 43,1 ), e a menor, na região Norte ( 32,4 ).
As demais medianas foram : 37,0 na região Nordeste, 33,7 na região Sudeste e 35,5 na região Centro-Oeste.
A nota Máxima do Brasil, como um todo, foi 91,6, sendo obtida por, pelo menos, um aluno na região Sul.
As demais notas máximas foram : 65,4 na região Norte, 79,2 na região Nordeste, 90,1 na região Sudeste e 81,0 na região Centro-Oeste.
A nota Mínima foi zero nas regiões Nordeste e Sudeste, 3,7 nas regiões Sul e Centro-Oeste e 7,4 na região Norte.
Observa -se que existe diferença estatisticamente significativa da maior Média das notas do Componente de Conhecimento Específico da região Sul ( 43,7 ) em relação às médias das demais regiões.
Assim como os Gráficos 6.1 e 6.2, o Gráfico 6.3, apresentado a seguir, permite uma avaliação do desempenho de concluintes de Engenharia de Controle e Automação em relação ao Componente de Conhecimento Específico com um histograma da distribuição das notas correspondentes.
Esta também é uma distribuição unimodal, e o grupo modal é o ( 30 ; 40 ], seguido de perto do intervalo ( 20 ; 30 ], abaixo do grupo modal da prova como um todo e do grupo modal para a Formação Geral.
A Tabela 6.6 apresenta uma comparação dos resultados em relação à Categoria Administrativa, à Organização Acadêmica e a Modalidade de Ensino, agora levando em conta o desempenho de alunos do Componente de Conhecimento Específico da prova.
No que se refere à Organização Acadêmica, a maior Média foi das Universidades ( 40,4 ), seguida da Média dos CEFET/IFET ( 39,3 ), Centros Universitários ( 33,2 ) e Faculdades ( 31,2 ).
O maior Desvio padrão, e acima do valor para o Brasil, como um todo, foi o das Universidades ( 15,0 ).
As Universidades também obtiveram a maior nota Máxima ( 91,6 ).
Os Centros Universitários obtiveram nota Máxima igual a 90,1, os CEFET/IFET igual a 88,5 e as Faculdades obtiveram nota Máxima igual a 81,0.
As medianas foram : 39,7 nas Universidades, a maior delas, 37,5 nos CEFET/IFET, 32,1 nos Centros Universitários e 29,6 nas Faculdades, a menor.
A nota Mínima foi zero nas Universidades e nos CEFET/IFET.
Para as demais Organizações Acadêmicas a nota Mínima foi 3,7.
Observa -se que existem diferenças estatisticamente significativas ao nível de 95 \% no Componente de Conhecimento Específico entre as notas das Organizações Acadêmicas, exceto entre as Universidades e os CEFET/IFET.
Quanto à Categoria Administrativa, observa -se um comportamento semelhante àquele da parte de Formação Geral e à prova como um todo, ou seja, existe diferença estatisticamente significativa entre as médias das IES Públicas ( 42,4 ) e IES Privadas ( 32,8 ).
Neste caso também, a maior Média foi obtida por alunos de IES Públicas de ensino.
No que se refere à Modalidade de Ensino, como mencionado anteriormente, só houve estudantes de Educação Presencial, cuja Média foi 37,0. 6.2 ANÁLISE DAS QUESTÕES OBJETIVAS Esta seção apresenta estatísticas selecionadas e histogramas das Questões Objetivas de Formação Geral ( 6.2.1 ) e de Conhecimento Específico ( 6.2.2 ).
São também apresentadas e comparadas as médias das subpopulações caracterizadas por Grande Região. 6.2.1 Componente de Formação Geral A Tabela 6.7 apresenta as Estatísticas Básicas relativas às oito questões objetivas do componente da prova que abrange a Formação Geral dos estudantes.
A Média do Brasil foi 59,3.
A menor Média foi encontrada na região Norte ( 55,5 ), e a maior, na região Sul ( 64,9 ).
As demais médias foram : 62,2 na região Nordeste, 57,9 na região Sudeste e 61,0 na região Centro-Oeste.
O Desvio padrão do Brasil foi 21,9, sendo o maior Desvio padrão encontrado na região Sudeste ( 22,0 ), e o menor, na região Norte ( 19,6 ).
Os demais desvios foram : 21,8 na região Nordeste e 20,8 nas regiões Sul e Centro-Oeste.
As medianas do Brasil como um todo e de quatro das cinco regiões foi 62,5 - a única região com Mediana diferente foi a Norte, com 50,0.
A nota Máxima 100,0 foi alcançada em todas as grandes Regiões.
As notas mínimas ( 0,0 ) ocorreram nas regiões Nordeste, Sudeste e Sul.
As regiões Norte e Centro-Oeste tiveram nota Mínima igual a 12,5.
A Tabela 6.8 apresenta o Índice de Facilidade e o Índice de Discriminação ( ponto-bisserial ) para cada uma das questões objetivas do Componente de Formação Geral.
Quanto ao Índice de Facilidade, foram usadas as seguintes cores para diferenciar o nível de dificuldade da questão : A Tabela 6.9 apresenta a distribuição das questões considerando -se simultaneamente a classificação dos dois índices.
As sete questões com Índice de Discriminação Muito bom figuraram entre quatro níveis de dificuldade Muito fácil, Fácil, Médio e Difícil : uma na categoria Muito fácil ( questão 6 ) do Índice de Facilidade, duas classificadas na categoria Fácil ( questões 4 e 5 ), duas na categoria Médio ( questões 3 e 8 ) e, por fim, duas na categoria Difícil ( questões 1 e 7 ).
Em particular, as questões 1, 4 e 5 foram as que apresentaram o maior poder discriminatório, as três com índice 0,51, e foram consideradas uma Difícil e duas fáceis em termos de facilidade, com proporções de acertos, respectivamente, de 0,40, 0,73 e 0,72.
O máximo de acertos foi alcançado pela questão 6 com um Índice de Facilidade de 0,87.
A questão de número 2 apresentou um Índice de Facilidade de 0,55 ( Médio ), ou seja, pouco mais da metade dos estudantes conseguiu resolvê -la, dentro do universo de participantes.
Seu Índice de Discriminação foi Bom ( 0,39 ).
O Gráfico 6.4, para exemplificar, analisa o comportamento da questão de número 4 de Formação Geral.
Trata -se de uma questão considerada Fácil, com relação à facilidade e uma das que obteve o maior Índice de Discriminação dessa parte da prova ( 0,51 ), juntamente com as questões 1 e 5.
Neste gráfico, cada uma das cinco curvas representa o percentual de respostas em determinada alternativa da questão, em função do número de acertos dos estudantes nessa parte da prova ( Formação Geral/Múltipla Escolha ), antes de possíveis eliminações pelo critério do ponto-bisserial.
A curva em verde corresponde à alternativa B, a alternativa correta para esta questão.
Observa -se que entre os estudantes com menor número de acertos, nessa parte do exame, a situação mais frequente foi a escolha da alternativa E ( em vermelho ), incorreta.
Por exemplo, entre os estudantes que acertaram três questões, 19,2 \% escolheram a alternativa E, 49,1 \% escolheram a alternativa B ( em verde ), a alternativa correta, 17,9 \% escolheram a alternativa D ( em roxo ), 8,5 \% escolheram a alternativa A ( em azul ) e 5,1 \% a C ( em preto ).
Entre os que acertaram três respostas entre as questões de múltipla escolha de Formação Geral, 0,0 \% deixou a questão em branco e 0,2 \% marcou mais de uma alternativa, invalidando a questão.
À medida em que o número de acertos aumenta, indicando desempenho melhor nessa parte da prova, aumenta concomitantemente a proporção de estudantes que selecionaram a alternativa correta B, atingindo 100 \% para os estudantes com 8 acertos.
Essa análise permite verificar como a questão discriminou os grupos de desempenho, justificando o alto índice obtido na questão.
Cumpre notar que não é possível inferir deste gráfico nem o índice de facilidade ( que seria uma média da proporção ponderada pela quantidade de alunos com cada uma das notas ), nem o índice de discriminação ponto-bisserial, por razão equivalente.
No caso extremo no qual a grande concentração dos acertos dos alunos fosse abaixo de 4, o índice de facilidade seria obrigatoriamente abaixo de 40 \% ( neste exemplo ).
Caso a concentração fosse em 6 acertos ou mais, o índice seria obrigatoriamente acima de 60 \%.
Os gráficos relativos às demais questões de Formação Geral constam do Anexo I. 6.2.2 Componente de Conhecimento Específico A Tabela 6.10 apresenta as Estatísticas Básicas em relação às questões objetivas do Componente de Conhecimento Específico da prova de Engenharia de Controle e Automação por Grande Região.
Como será discutido ainda nessa seção, quatro questões objetivas do Componente de Conhecimentos Específicos não foram usadas no cômputo das notas por terem sido descartadas pelo critério do ponto-bisserial.
Assim, as notas foram calculadas com base em 23 das 27 questões objetivas de Conhecimentos Específicos.
A Média do Brasil deste componente foi de 41,6.
A menor Média foi observada na região Norte ( 37,4 ), e a maior, na região Sul ( 48,2 ).
O Desvio padrão de todo o Brasil foi 15,2, sendo o menor Desvio padrão encontrado na região Norte ( 12,5 ), e o maior, na região Sul ( 15,9 ).
A Mediana de todo o Brasil foi 39,1, igual à das regiões Sudeste e Centro-Oeste.
A região Norte teve a Mediana menor ( 34,8 ) e a região Sul a maior ( 47,8 ).
A região Nordeste teve a Mediana 43,5.
A nota Máxima da prova foi 91,3, obtida nas questões objetivas do Componente de Conhecimento Específico, por, pelo menos, um aluno das regiões Sudeste e Sul.
A nota Mínima foi zero em duas das cinco regiões, Nordeste e Sudeste.
As regiões Sudeste e Centro-Oeste tiveram a nota Mínima, igual a 4,3 e a região Norte teve a nota Mínima maior, igual a 8,7.
A diferença entre a nota Média mais alta obtida na região Sul ( 48,2 ) e a nota Média das demais regiões é estatisticamente significativa.
A Tabela 6.11 apresenta os Índices de Facilidade e Discriminação ( ponto-bisserial ) das questões objetivas do Componente de Conhecimento Específico para os estudantes de Engenharia de Controle e Automação.
Para facilitar a diferenciação das questões, usaremos as mesmas cores da Tabela 6.8 para as diferentes classificações dos Índices de Facilidade e de Discriminação.
A partir do Índice de Facilidade obtido, pode -se concluir que mais de metade das questões objetivas da prova foi considerada, pelo menos, Difícil : das 27 questões, 16 foram classificadas como Difícil ou como Muito difícil.
Oito questões foram classificadas como Médio, e outras três consideradas como Fácil.
Nenhuma questão foi classificada como Muito fácil.
Já quanto aos índices de discriminação das questões objetivas do Componente de Conhecimento Específico da prova, tem -se como resultado a seguinte classificação : 11 das 27 questões foram consideradas como boas, enquanto três delas tiveram Índice de Discriminação Muito bom.
Assim, para um pouco mais da metade das questões – 14 em 27 – os Índices de Discriminação foram Bom ou Muito bom.
Dentre as demais, nove delas foram classificadas como Médio e outras quatro como Fraco, sendo 13, por conseguinte, a quantidade de questões nos dois patamares mais baixos de discriminação.
Constata -se, assim, que a prova – no que se refere ao Componente de Conhecimento Específico – possuía boa capacidade de discriminar entre aqueles que dominavam ou não o conteúdo.
A Tabela 6.12 apresenta a distribuição das questões considerando -se simultaneamente a classificação dos dois índices.
Dentre as questões que alcançaram os maiores índices de discriminação, classificadas com Índice de Discriminação Muito bom, uma delas, a de número 28, de índice mais elevado ( 0,42 ), foi classificada como Médio quanto ao Índice de Facilidade, com 48 \% dos estudantes marcando a opção correta.
A questão de número 33 foi a mais difícil dentre as 27 questões específicas, com baixo Índice de Facilidade, apenas 13 \% de acertos.
Essa questão apresentou poder discriminatório muito baixo, 0,19, o que comprova ter sido a mais difícil para os estudantes.
Destaca -se, também, a questão 20, com Índice de Facilidade 0,17, o que, em termos percentuais, corresponde a 17 \% de estudantes que responderam acertadamente.
Já 0,16 foi o seu de Índice de Discriminação.
Tais questões foram, portanto, pelo critério ponto-bisserial, consideradas inadequadas.
Por isso, as questões 11 e 17, além das questões 20 e 33 foram eliminadas do cômputo da nota final.
A título de exemplo das análises do comportamento das questões objetivas, o Gráfico 6.5 analisa a questão 28 do Componente de Conhecimento Específico.
Esta questão foi considerada pelas respostas dos estudantes avaliados na prova como uma questão com nível de facilidade Médio, 0,48, ou seja, 48 \% dos estudantes assinalaram acertadamente a opção E, correspondente ao gabarito.
Como já comentado, seu Índice de Discriminação foi igual a 0,42, classificado como Muito bom, o maior valor de discriminação.
Neste gráfico, cada uma das cinco curvas representa o percentual de respostas em determinada alternativa da questão 28, em função do número de acertos dos estudantes nessa parte da prova, antes de possíveis eliminações de questões pelo critério do ponto-bisserial.
A alternativa correta E, representada no gráfico pela curva em vermelho, foi escolhida em maiores proporções pelos alunos com desempenho melhor nessa parte da prova.
Já as alternativas incorretas, também denominadas distratores, foram selecionadas, principalmente, por aqueles com notas mais baixas.
Observa -se que a soma não é 100 \% por causa das questões não respondidas ou com mais de uma opção marcada.
Aqueles com nota zero, na sua quase totalidade deixaram esta questão em branco ou marcaram mais de uma alternativa, comportamento considerado inválido.
A proporção de alunos que selecionou a resposta correta E aumenta gradativamente, chegando a atingir 100 \% para 21 acertos ou mais, enquanto a proporção dos que escolheram alternativas incorretas decai, a partir de três ou quatro acertos, como função do número de acertos nessa parte da prova.
Os gráficos relativos às demais questões do Conhecimento Específico constam do Anexo I. 6.3 ANÁLISE DAS QUESTÕES DISCURSIVAS Esta seção apresenta estatísticas selecionadas e histogramas das Questões Discursivas de Formação Geral ( 6.3.1 ) e Conhecimento Específico ( 6.3.2 ).
São também apresentadas e comparadas as médias de subpopulações caracterizadas por Grande Região. 6.3.1 Componente de Formação Geral As análises dos resultados de desempenho dos estudantes de Engenharia de Controle e Automação nas duas questões discursivas relativas à Formação Geral encontram -se na Tabela 6.13 e no Gráfico 6.6.
Na Tabela 6.13, observa -se que a nota Média nesse conjunto de questões ficou abaixo da obtida nas objetivas.
Os estudantes, de todo o Brasil, obtiveram, em Formação Geral, Média 59,3 nas questões objetivas e 54,5 nas questões discursivas.
Pode -se notar também que o Desvio padrão nesse conjunto de questões foi maior do obtido nas objetivas, 21,9 nas questões objetivas e 24,9 nas questões discursivas.
A maior Média foi obtida na região Norte ( 57,8 ), e a menor, na região Centro-Oeste ( 50,9 ).
A Mediana de todo o Brasil, neste componente, foi 61,0, a mesma sendo obtida na região Sudeste.
Nas regiões Nordeste e Centro-Oeste a Mediana foi menor ( 58,3 e 58,5, respectivamente ) e nas demais regiões foi maior que o valor para o Brasil : Norte ( 63,3 ) e Sul ( 63,0 ).
A nota Máxima ( 97,5 ) foi a mesma nas regiões Sudeste e Sul, sendo 96,5 na região Norte, 95,0 na região Nordeste e 95,5 na região Centro-Oeste.
A nota Mínima ( 0,0 ) foi a mesma em todas as regiões do Brasil.
O Gráfico 6.6 representa a distribuição das notas nas questões discursivas do Componente de Formação Geral.
A moda desta distribuição ocorre no intervalo ( 60 ; 70 ], com frequência um pouco maior do que a do intervalo ( 70 ; 80 ].
Destaca -se também o intervalo [ 0 ; 10 ] com distribuição maior que 10 \% do total de notas, sendo que no intervalo [ 0 ; 10 ] incluem -se além da nota zero, a frequência de alunos que deixaram este tipo de questão em branco.
Na sequência, os resultados verificados para cada uma das questões discursivas de Formação Geral serão apresentados, estabelecendo -se relações com os conteúdos abordados em cada uma delas.
Os comentários da Banca de docentes corretores, a respeito do observado na correção das respostas dos estudantes, suas impressões e conclusões serão apresentados junto à análise de cada questão.
Cumpre esclarecer que, tendo em vista que as questões discursivas de Formação Geral são padronizadas, ou seja, constam de todas as provas, os comentários da Banca são os mesmos para todas as carreiras acadêmicas, sendo direcionados a todos os estudantes que participaram do Enade/2017.
A seguir, serão analisados os desempenhos da Área de Engenharia de Controle e Automação nas duas questões discursivas de Formação Geral do Enade/2017, comparando -se os resultados obtidos com comentários para cada questão. 6.3.1.1 Análise de Conteúdo da Questão Discursiva 1 do Componente de Formação Geral Os dados de Engenharia de Controle e Automação, obtidos a partir das respostas à questão 1, encontram -se na Tabela 6.14 e no Gráfico 6.7.
Nessa questão – de melhor desempenho dentre as duas de Formação Geral – os alunos, de todo o Brasil, tiveram Média 55,6.
A maior Média para a questão 1 foi obtida na região Norte ( 60,3 ), e a menor, na região Centro-Oeste ( 52,0 ).
Quanto à variabilidade das notas, o Desvio padrão de todo o Brasil foi 28,1.
O menor Desvio padrão foi obtido na região Norte ( 26,0 ), e o maior Desvio padrão foi obtido na região Nordeste ( 29,6 ).
As medianas de todas as regiões e do Brasil como um todo, foram iguais ( 60,0 ), com exceção da região Norte, com Mediana 65,0.
As notas máximas e mínimas da questão discursiva 1 foram, sem exceção, as mesmas para todas as regiões do Brasil, respectivamente, 100,0 e 0,0.
O Gráfico 6.7 mostra a distribuição das notas na questão discursiva 1 do Componente de Formação Geral.
Dentre os que responderam à questão a moda foi o intervalo ( 70 ; 80 ], com aproximadamente 23 \% do total de estudantes.
Observa -se um outro máximo local correspondente aos alunos que deixaram a resposta a essa questão em branco, em torno de 11 \% dos participantes. 6.3.1.2 Comentários sobre a correção de Conteúdo das respostas à Questão Discursiva 1 O enunciado era claro e indicava um tema de extrema relevância, qual seja, a incidência de doença grave integrante do rol das DST ( Doenças Sexualmente Transmissíveis ), que encaminha para a análise do direito fundamental à saúde, que tem assento constitucional ( CRFB, art. 6º. caput : São direitos sociais a educação, a saúde, a alimentação, o trabalho, a moradia, o transporte, o lazer, a segurança, a previdência social, a proteção à maternidade e à infância, a assistência aos desamparados, na forma desta Constituição ).
A questão focalizava, por meio de três textos, um tema social extremamente relevante, qual seja, a relação entre homem e mulher nas suas interações íntimas, no caso representadas pela vida sexual, e o reflexo das relações de poder veiculadas pela sociedade, que apontam, no mais das vezes, o predomínio do homem.
Tal posição de superioridade é Um dos textos demandava uma reflexão sobre a realidade dos casais que convivem em relações estáveis ou fugazes e não têm o controle dos efeitos das relações sexuais, que abrangem a gravidez não planejada e, em decorrência do não cuidado, a inserção de doenças transmitidas sexualmente.
Esse panorama é refletido na pouca influência das mulheres sobre os homens quanto à proteção, normalmente representada pelo uso de preservativos nas relações intimas, e após a constatação da doença, no encaminhamento das pessoas para tratamento.
Como efeito colateral, está a transmissão de doenças para os fetos.
Outro texto trazia à tona o maior cuidado que as mulheres têm, em geral, com sua saúde, notadamente no momento em que se descobre grávida e que redunda na realização de exames obrigatórios de pré-natal, que podem facilmente detectar doenças e, caso encontradas, ter o seu tratamento encaminhado.
A par disso, destaca -se que a crise de recursos financeiros agravou a situação, que já recebe influxos sociais negativos, como a dificuldade do homem em utilizar preservativos, realizar exames e procurar tratamentos quando adoece.
Como muito divulgado pelas mídias, a crise financeira do país repercutiu diretamente na prestação de serviços de saúde em todas as esferas : distrital, federal, estadual e municipal.
Assim, agravou -se o problema da falta de medicamentos essenciais para o tratamento da sífilis e, certamente, de outras doenças.
Como se verifica, a temática indicada no enunciado, representada pelos textos, pertence à realidade conhecida amplamente pelos estudantes, pois também é divulgada por todos os meios de comunicação.
Trata -se de tema geral que, no entanto, normalmente, não se inclui nas grades curriculares de cursos superiores que não sejam da área de Saúde.
Apesar disso, a temática é de conhecimento geral e relacionada aos direitos humanos, esses sim, ponto de integração com as diretrizes curriculares.
A linguagem utilizada era clara, indicando a existência dos fatos vinculados a questões relevantes e à saúde pública.
No entanto, como os textos trazem uma ampla gama de possibilidades de interpretação, ocorreram algumas compreensões diversas das esperadas.
Além disso, ressalta -se que a amplitude dos textos disponibilizados no enunciado fornece elementos que poderiam ser simplesmente reproduzidos pelos concluintes nas suas respostas.
O enunciado da questão encaminha a resposta para dois aspectos.
O estudante deveria abordar conteúdo pertinente à vulnerabilidade das mulheres às DST e o papel social do homem em relação à prevenção dessas doenças.
Para responder a esses aspectos, um justificada pelas convenções arraigadas que acarretam a educação dos homens e mulheres por padrões patriarcais, que redundam em posições machistas.
Um dos textos demandava uma reflexão sobre a realidade dos casais que convivem em relações estáveis ou fugazes e não têm o controle dos efeitos das relações sexuais, que abrangem a gravidez não planejada e, em decorrência do não cuidado, a inserção de doenças transmitidas sexualmente.
Esse panorama é refletido na pouca influência das mulheres sobre os homens quanto à proteção, normalmente representada pelo uso de preservativos nas relações intimas, e após a constatação da doença, no encaminhamento das pessoas para tratamento.
Como efeito colateral, está a transmissão de doenças para os fetos.
Outro texto trazia à tona o maior cuidado que as mulheres têm, em geral, com sua saúde, notadamente no momento em que se descobre grávida e que redunda na realização de exames obrigatórios de pré-natal, que podem facilmente detectar doenças e, caso encontradas, ter o seu tratamento encaminhado.
A par disso, destaca -se que a crise de recursos financeiros agravou a situação, que já recebe influxos sociais negativos, como a dificuldade do homem em utilizar preservativos, realizar exames e procurar tratamentos quando adoece.
Como muito divulgado pelas mídias, a crise financeira do país repercutiu diretamente na prestação de serviços de saúde em todas as esferas : distrital, federal, estadual e municipal.
Assim, agravou -se o problema da falta de medicamentos essenciais para o tratamento da sífilis e, certamente, de outras doenças.
Como se verifica, a temática indicada no enunciado, representada pelos textos, pertence à realidade conhecida amplamente pelos estudantes, pois também é divulgada por todos os meios de comunicação.
Trata -se de tema geral que, no entanto, normalmente, não se inclui nas grades curriculares de cursos superiores que não sejam da área de Saúde.
Apesar disso, a temática é de conhecimento geral e relacionada aos direitos humanos, esses sim, ponto de integração com as diretrizes curriculares.
A linguagem utilizada era clara, indicando a existência dos fatos vinculados a questões relevantes e à saúde pública.
No entanto, como os textos trazem uma ampla gama de possibilidades de interpretação, ocorreram algumas compreensões diversas das esperadas.
Além disso, ressalta -se que a amplitude dos textos disponibilizados no enunciado fornece elementos que poderiam ser simplesmente reproduzidos pelos concluintes nas suas respostas.
O enunciado da questão encaminha a resposta para dois aspectos.
O estudante deveria abordar conteúdo pertinente à vulnerabilidade das mulheres às DST e o papel social do homem em relação à prevenção dessas doenças.
Para responder a esses aspectos, um número expressivo de concluintes utilizou somente os elementos dos textos motivadores, sem nada acrescentar.
A resposta deveria conter também uma segunda parte, a indicação de duas ações voltadas para o público masculino, a serem adotadas no âmbito das políticas públicas de saúde ou de educação, para a redução do problema.
Verificou -se que o desempenho dos estudantes nessa questão foi de médio para bom.
Ao se analisarem as respostas, constata -se que muitos foram os que apresentaram os tópicos presentes no padrão de resposta, amplo o suficiente para albergar as respostas apresentadas pelos concluintes.
Tal resultado pode ser justificado, tendo em vista tratar -se de tema bastante difundido nas diversas mídias, por ter sido apresentado de forma ampla pelos textos motivadores presentes no enunciado.
Por outro lado, o debate no ambiente universitário sobre as relações de poder entre homens e mulheres tende a ser marcante pelo movimento de empoderamento feminino, liderado por organizações sociais com ampla atuação nos estabelecimentos de nível superior.
De forma recorrente, as relações patriarcais entre homens e mulheres foram vinculadas ao modelo machista vigente.
Essas condições, segundo os concluintes, explicam a incidência da doença e a situação de vulnerabilidade das mulheres, bem como a função social dos homens.
Outras respostas apontaram o poder público, nas áreas de saúde e/ou educação, como órgãos responsáveis pela solução do problema.
Para os que se fixaram na área da saúde, foram frequentes a indicação da necessidade de mais medicamentos, vacinas, educação sexual, médicos e postos de atendimentos.
Alguns, no entanto, apenas indicaram ações individuais para o homem, e outros, ações individuais para as mulheres.
Também figuraram respostas que apresentavam ações da igreja como solução, além de outras que faziam referências a campanhas no seio das famílias. 6.3.1.3 Análise de Conteúdo da Questão Discursiva 2 do Componente de Formação Geral A Tabela 6.15 mostra que o desempenho médio dos estudantes na questão discursiva 2 ( média 53,1 ) foi inferior ao obtido na questão discursiva 1 ( média 55,6 ).
As regiões Norte e Sul foram aquelas cuja Média, também nessa questão, foi maior ( 56,3 em ambas ), e a de menor Média foi a da região Nordeste ( 48,6 ).
Quanto à variabilidade das notas, o Desvio padrão de todo o Brasil foi 33,5, superior ao obtido na questão discursiva 1 ( 28,1 ).
O maior desvio nessa questão foi obtido na região Nordeste ( 36,0 ), enquanto o menor foi obtido na região Sul ( 32,2 ).
A Mediana de todo o Brasil foi 70,0, a mesma de todas as regiões, com exceção das regiões Nordeste e Centro-Oeste que tiveram Mediana menor, 50,0.
As notas máximas ( 100,0 ) e as notas mínimas ( 0,0 ) também foram as mesmas em todas as regiões, sem exceção.
O Gráfico 6.8 mostra a distribuição das notas na questão discursiva 2 do Componente de Formação Geral.
Observa -se que a maior frequência corresponde aos alunos que obtiveram nota dentro do intervalo ( 70 ; 80 ].
Destacam -se, também, os alunos que deixaram a questão em branco e os alunos que obtiveram nota dentro do intervalo ( 90 ; 100 ], caracterizando -se como máximos locais. 6.3.1.4 Comentários sobre a correção de Conteúdo das respostas à Questão Discursiva 2 O enunciado era claro e abordava um tema de extrema relevância : a questão de gênero.
Trata -se de uma proposta polêmica, pois abrange a análise da inserção social de minoria vinculada à questão de gênero.
O enunciado solicitava que o concluinte discorresse sobre a importância do nome para as pessoas transgêneras e propusesse uma medida, no âmbito das políticas públicas, para facilitar o acesso dessas pessoas à cidadania.
A respeito da primeira parte, a importância do nome, esperava -se que o estudante abordasse as dificuldades que o uso do nome civil ( nome de registro, nome de batismo, nome da certidão de nascimento, por exemplo ) pode acarretar para as pessoas transgêneras, além de demonstrar conhecimento sobre a importância do nome social, direito já reconhecido através, principalmente, de decisões judiciais e atos administrativos, como decretos e portarias que dispõem sobre o uso do nome social.
Em atendimento à segunda parte do comando, o participante deveria propor uma política pública associada ao direito à cidadania das pessoas transgêneras.
Sobre esse aspecto, esperava -se a compreensão do significado de política pública, ou seja, propostas de âmbito governamental, e não, iniciativas de Organizações Sociais, igrejas ou grupos sociais de qualquer tipo.
O nascimento das pessoas, no Brasil, é objeto de inscrição em cartórios de Registro Civil das Pessoas Naturais, órgão responsável por estabelecer a certificação pública dos dados de todos os indivíduos que não somente nascem, mas casam, eventualmente se divorciam e vêm a óbito.
Daí o nome civil.
Para flexibilizar a correção, outras nomenclaturas foram consideradas, a saber : identificação de nascimento, nome de batismo, nome de registro, nome próprio, nome verdadeiro, que diferenciam o nome civil daquele que o indivíduo deseja utilizar para traduzir a sua opção pessoal de gênero, que se convencionou denominar nome social.
Quanto ao entendimento do termo transgênero, deve -se destacar que o movimento de mudança do gênero é contemporâneo e não está necessariamente vinculado à mudança de sexo.
Existem várias hipóteses em jogo.
Uma é a que envolve cirurgias de mudança de sexo ( a menina extirpa os seus seios e inclui órgão genital por intervenção cirúrgica ; o menino inclui seios e extirpa o seu órgão genital e inclui o órgão feminino também por intervenção cirúrgica ).
Há outro movimento que envolve apenas as vestimentas e o linguajar corporal, além dos relacionamentos ( as meninas assumem o trajar, o andar, os gostos dos meninos e mantêm relacionamentos afetivos com meninas ; os meninos assumem o trajar, o andar, os gostos das meninas e mantêm relacionamentos afetivos com meninos ).
Em relação à discussão da primeira solicitação do comando da questão, destaca -se que todos têm nome civil, que foi escolhido pelos seus pais ou responsáveis e, na falta deles, por autorização do Juiz de Direito, pautado no sexo detectado no nascimento.
Já o nome social é aquele que indica a preferência do indivíduo na sua identificação nas relações sociais e que destoa do seu gênero de registro civil.
Ressalte -se que o Brasil não possui ainda legislação específica sobre pessoas transgêneras.
Observa -se um relativo progresso no reconhecimento de direitos, através, principalmente, de decisões judiciais e atos administrativos, como decretos e portarias que dispõem sobre o uso do nome social.
Sem uma lei que defina os procedimentos da alteração dos documentos para pessoas transexuais, essa parcela da população LGBT ( lésbicas, gays, bissexuais, transexuais e transgêneros ) é obrigada a procurar na Justiça o reconhecimento de sua identidade, em processos que podem ser longos e que dependem do olhar de determinados, profissionais, como médicos, psicólogos, assistentes sociais, bem como do entendimento dos advogados, defensores públicos, juízes e do Judiciário, como um todo.
Atualmente, tramita na Comissão de Direitos Humanos da Câmara dos Deputados o Projeto de Lei João Nery ( 5002/2013 ), dos deputados Jean Willys ( PSOL-RJ ) e Erika Kokay ( PT-DF ), que determina que o reconhecimento da identidade de gênero seja um direito do cidadão.
O projeto recebeu o nome do primeiro trans-homem operado no Brasil.
Grande parte dos respondentes considerou a incongruência entre a imagem corporal, representada pelo gênero assumido pelos trans, e o nome que apresenta nos seus documentos como causa de sérios constrangimentos e embaraços e as dificuldades para que esses sujeitos acessem direitos básicos, como saúde, educação e empregabilidade em uma sociedade heteronormativa e preconceituosa.
Muitos demonstraram reconhecer que a identidade e expressão de gênero, através do nome social, pressupõe que sejam respeitadas a dignidade da pessoa humana, a intimidade, a liberdade e o direito à cidadania plena.
No entanto, também foram encontradas, em quantidades significativas, respostas que demonstravam a existência de uma interpretação social conservadora que não admite a mudança de sexo nem do nome social.
Existe um acentuado preconceito contra todos os movimentos que buscam identificação social diversa da civil.
Isso inclui os variados grupos : gays, travestis, transgêneros, dentre outros.
Já em relação ao segundo tópico a ser abordado pelo concluinte, é notório que as pessoas trans encontram -se destituídas de direitos no campo das políticas públicas, como saúde, habitação, emprego, educação, justiça, dentre outros.
Nesse sentido, uma grande diversidade de medidas, no âmbito das políticas públicas, poderia ser proposta.
Ações voltadas para a capacitação de profissionais da educação, saúde, segurança pública, dentre outras áreas, para melhor atender às demandas e especificidades desse público alvo seriam necessárias.
Ainda em relação à educação, seriam também necessárias ações que contribuíssem para diminuir a evasão, pois o ambiente escolar é difícil para uma pessoa transgênera, por conta das situações discriminatórias e preconceituosas vindas dos colegas de classe, e mesmo das/dos professores e/ou de outros profissionais da área, não podem ser mais um complicador nesse processo.
Ainda no que tange ao assunto, percebe -se que, ao serem excluídas dos bancos escolares, as oportunidades no mercado de trabalho formal praticamente inexistem.
Um outro aspecto que se evidencia no cotidiano desse segmento da sociedade diz respeito às situações de violência sofridas pelas pessoas transgêneras.
Nesse sentido, algumas propostas foram elencadas para minimizar o problema e combater a violência.
Políticas envolvendo melhor formação de funcionários das polícias civil e militar, por exemplo, somariam esforços na qualificação das abordagens policiais e dos atendimentos prestados pelas forças de segurança, pois o respeito à orientação sexual e à identidade e expressão de gênero é, antes de tudo, o respeito à dignidade, à cidadania e à própria democracia.
Nesse sentido, revela -se essencial que, ao sofrerem agressões, sejam elas físicas, psicológicas ou morais no espaço público ou domiciliar, as pessoas trans possam ser atendidas, com respeito inclusive ao nome social, nas delegacias comuns e, inclusive, nas Delegacias Especiais de Atendimento à Mulher ( DEAM ), pois esse é o órgão institucional que possivelmente detém as qualificações necessárias para o atendimento ao gênero feminino ao qual as travestis e mulheres trans efetivamente pertencem.
Também ocorre uma latente inacessibilidade das pessoas trans à assistência jurídica.
Nesse sentido, é necessário implementar projetos de capacitação para profissionais do poder judiciário, como um todo, de modo a prestarem a assistência jurídica gratuita na defesa intransigente dos direitos sociais, civis, trabalhistas e políticos dessa população.
Dessa maneira, a Justiça agiria de forma mais acessível e sensível para o segmento, que terá suas especificidades observadas a começar pelo devido respeito do direito ao uso do nome social.
Nesse âmbito, o mais comum foi propor a adoção de medidas punitivas para quem viola o direito à autodeterminação de gênero e a elaboração de leis que garantam a mudança do nome e assegurem outros direitos para as pessoas trans.
Propostas envolvendo a formação de profissionais, como as mencionadas aqui, e de conscientização da população foram as mais frequentes.
Muitas outras enfatizavam o direito ao uso respeitado do nome social e a garantia de alteração do registro civil, independentemente da cirurgia de redesignação sexual ou transgenitalização.
Muitos questionaram o fato de não bastar a autodeclaração no processo de alteração de nome, enfatizando que o nome social é uma questão de princípio de dignidade e cidadania.
Apesar de pouco mencionada nas respostas corrigidas, políticas envolvendo a promoção da saúde são muito necessárias.
Cabe aos gestores de saúde ações direcionadas, não somente à prevenção e tratamento do HIV/Aids, mas também políticas de saúde que contemplem as especificidades dessa população, tais como o Ministério da Saúde, especialmente o SUS ( Sistema Único de Saúde ), incluir, de fato, o atendimento para pessoas trans para acompanhamento da terapia hormonal, que muitas demandam, respeitando -se o direito de utilização do nome social.
De modo geral, são muitas as políticas públicas a que precisam ser implantadas.
Urge a necessidade de uma agenda mais compromissada por parte dos gestores do nosso país com essa comunidade, que ainda é muito discriminada.
Logo, tais políticas devem contemplar suas necessidades básicas como o direito de acesso aos estudos, à capacitação para o mercado de trabalho, à profissionalização, ao acesso a bens e serviços de qualidade em saúde, habitação, segurança, cultura e assistência social e, certamente, o respeito ao uso do nome social, o que também contribui para essa inclusão.
Assim, as opções apresentadas no padrão de resposta funcionaram como exemplos de políticas públicas possíveis, já que seria difícil apresentar uma lista suficientemente ampla para abranger a grande variedade de respostas encontradas.
O desempenho dos estudantes nessa questão foi médio.
As respostas, de modo geral, indicaram a necessidade de uma atualização por parte dos concluintes quanto a temas essenciais e, no caso da presente questão, a análise sobre aspectos fundamentais da questão de gênero.
Muitos, possivelmente motivados pelos textos apresentados no enunciado, restringiram -se a discorrer sobre preconceito de gênero ou mesmo sobre o preconceito de modo mais geral, sem responder ao que foi solicitado no comando da questão.
Outros apresentaram discursos religiosos ou ideológicos, também sem tratarem do que foi solicitado. 6.3.1.5 Análise de Língua Portuguesa das Questões Discursivas do Componente de Formação Geral Os dados de Engenharia de Controle e Automação, obtidos a partir das respostas às questões discursivas do Componente de Formação Geral, no que tange à Língua Portuguesa, encontram -se na Tabela 6.16 e no Gráfico 6.9.
Nesse aspecto, os alunos, de todo o Brasil, obtiveram Média 55,4.
A maior Média com respeito à Língua Portuguesa foi obtida na região Sul ( 56,5 ), e a menor, na região Nordeste ( 51,7 ).
Quanto à variabilidade das notas, o Desvio padrão de todo o Brasil foi 23,0.
O menor Desvio padrão foi obtido na região Norte ( 18,7 ) e o maior Desvio padrão foi obtido na região Centro-Oeste ( 25,7 ).
A Mediana das notas de Língua Portuguesa foi 62,5 para o Brasil como um todo, a mesma obtida em duas regiões Sudeste e Sul.
Nas demais regiões a Mediana foi 60,0.
A nota Máxima para todo o Brasil foi de 95,0, com, pelo menos, um aluno tirando essa nota nas regiões Sudeste e Sul.
Na região Norte a nota Máxima foi 82,5, na região Nordeste foi 87,5 e na Centro-Oeste foi 85,0.
Já a nota Mínima foi zero em todas as regiões do país.
O Gráfico 6.9 mostra a distribuição das notas de Língua Portuguesa do Componente de Formação Geral.
Observa -se que a maior frequência ( cerca de 28 \% ) corresponde aos alunos que obtiveram nota no intervalo ( 60 ; 70 ].
Destaca -se, também, os alunos que deixaram ambas as questões em branco representando quase 10 \% do total, caracterizando -se como máximo local. 6.3.1.6 Comentários sobre a correção das respostas de Formação Geral com respeito à Língua Portuguesa Os enunciados apresentados em 2017 não explicitaram a exigência de elaboração de um texto “ dissertativo ”, como nos anos anteriores.
O participante deveria fazer uma exposição de seus conhecimentos sobre cada um dos assuntos e estruturar seus textos de acordo com as características do registro formal adequado à situação comunicativa – avaliação de conhecimentos.
Essa configuração determina exigências quanto aos seguintes aspectos : adequação da seleção vocabular, desenvolvimento do conteúdo, organização lógica das ideias, estruturação sintática dos períodos, utilização de procedimentos de encadeamento textual e de referenciação, obediência às exigências morfossintáticas próprias da modalidade escrita da norma-padrão, respeito às regras ortográficas e às regras de acentuação gráfica.
O padrão de resposta utilizado na avaliação das questões 1 e 2 considerou os aspectos relevantes ao bom desempenho linguístico como competências distintas, de modo a permitir um mapeamento detalhado do domínio dos recursos disponíveis na Língua Portuguesa para a comunicação escrita formal.
Com base nesse objetivo, foram avaliados os seguintes aspectos : ( a ) estruturação textual condizente com o gênero solicitado e o modo de organização textual expositivo adequado ao gênero – essa competência envolve : • estruturação sintática condizente com o padrão da modalidade escrita formal da língua portuguesa de modo a garantir a clareza necessária ; • distribuição do conteúdo do texto em parágrafos, de modo a garantir a sua organização temática ; • utilização de operadores discursivos que contribuam para a progressão temática do texto, estabelecendo relações lógicas entre as ideias apresentadas, tanto do ponto de vista intrafrasal, como do interfrasal ; • utilização de procedimentos de referenciação lexical e pronominal que permitam a retomada de referentes textuais ; • utilização de sinais de pontuação que contribuam para a organização lógica da frase e do texto ; • inteligibilidade relacionada ao atendimento das exigências de estruturação textual.
Espera -se, portanto, que o participante recorra a procedimentos linguístico-discursivos para organizar seu texto, permitindo o encadeamento lógico entre suas partes de forma a garantir a progressão e a coerência textuais.
Isso significa que os seguintes procedimentos foram penalizados, de acordo com o padrão de resposta proposto : • elaboração de frases fragmentadas que comprometam a estrutura lógico-gramatical do texto ; • sequência justaposta de ideias sem encaixamentos sintáticos, reproduzindo hábitos da oralidade ; • elaboração de frase com apenas oração subordinada, sem oração principal ; • emprego equivocado do conector ( preposição, conjunção, pronome relativo, alguns advérbios e locuções adverbiais ) comprometendo a expressão da relação lógica entre duas ideias, com prejuízo da clareza do texto ; • emprego do pronome relativo sem a preposição, quando obrigatória ; • repetição ou substituição inadequada de palavras sem o emprego dos recursos oferecidos pela língua ( pronome, advérbio, artigo, sinônimo ) ; • emprego inadequado dos pronomes relativos “ cujo(a) ” e “ onde ” ; • utilização inadequada dos sinais de pontuação, comprometendo a clareza textual ; • inteligibilidade relacionada ao atendimento das exigências de estruturação textual. ( b ) respeito às convenções ortográficas da norma-padrão da Língua Portuguesa – essa competência envolve o domínio das regras de acentuação gráfica e da grafia padrão das palavras ( com ausência de abreviaturas próprias da linguagem da internet ), de acordo com as convenções estabelecidas pela legislação em vigor e consubstanciadas no Vocabulário Ortográfico da Língua Portuguesa, editado pela Academia Brasileira de Letras. ( Vale notar que, nessa edição de 2017, não foi aceita a legislação anterior, no caso das regras relativas ao uso do hífen e da acentuação gráfica ).
Espera -se que o participante :  grafe corretamente as palavras ;  respeite as regras de acentuação gráfica ;  empregue maiúsculas em início de frase, em nomes próprios de pessoas, lugares ou instituições ;  grafe as siglas com letras maiúsculas ou apenas com a primeira letra maiúscula quando formarem uma palavra ;  evite abreviações como p/, vc, tb, pq, tá, né, usadas muitas vezes em escrita informal e na internet ;  obedeça às regras de separação de sílabas no final da linha. ( c ) domínio dos diferentes aspectos morfossintáticos próprios da modalidade escrita formal da norma-padrão da Língua Portuguesa – essa competência envolve : a concordância nominal, a concordância verbal, a regência nominal, a regência verbal, a flexão nominal, a flexão verbal, a correlação entre os tempos verbais, a colocação pronominal e a utilização de sinais de pontuação que contribuam para a organização lógica da frase e do texto.
Espera -se que o participante :  flexione o verbo para estabelecer concordância de número com o sujeito da frase ;  flexione o artigo, o adjetivo e o pronome para concordar em número e em gênero com o substantivo a que se referem ;  observe a regência nominal e a verbal, utilizando a preposição adequada depois de um substantivo, um verbo ou um adjetivo ;  empregue adequadamente o acento grave indicador da crase entre uma preposição e um artigo ( a+a ) ;  obedeça às regras de colocação pronominal ( próclise e ênclise ), distintas dos hábitos da oralidade ou da escrita informal ( exigência de próclise com termo atrator, não exigência de que o pronome oblíquo se ligue ao verbo auxiliar por meio de hífen ) ;  flexione adequadamente verbos, substantivos, adjetivos e pronomes no que diz respeito à expressão das categorias gramaticais ;  flexione os verbos para expressar a correlação de modo e tempo nas estruturas subordinadas.
Com base nesses critérios, foram considerados como desvios de caráter morfossintático, e não como desvios ortográficos, as alterações que envolvem mudança de classe gramatical ou de forma flexional do verbo :  eliminação da marca de infinitivo ( -r- ) e substituição por acento agudo ou ausência total de marca do infinitivo ;  confusão entre “ ão ” e “ am ” nas formas verbais ;  confusão entre “ há ” e “ a ” ;  uso de hífen para separar pronome átono – tanto uso indevido quanto omissão ( exemplo : “ esperasse ”, em lugar de “ espera -se ” ; “ falar -mos ”, no lugar de “ falarmos ” ;  “ esta ” ( no lugar de “ está ” ) ; “ mais ” ( no lugar de “ mas ” ) ; “ e ” ( no lugar de “ é ” ).  verbos “ ter ” e “ vir ” que, na terceira pessoa do plural, não apresentarem o acento circunflexo, serão penalizados como desvio de concordância, em aspectos morfossintáticos. ( d ) seleção vocabular adequada à modalidade escrita formal da Língua Portuguesa, exigida pela situação comunicativa – essa competência envolve a precisão na seleção/utilização do vocabulário relacionado à temática solicitada pela questão ; a ausência de marcas da oralidade, como termos de sentido muito genérico ( “ coisa ”, “ negócio ”, “ você ” ) e termos de registros mais informais ( como gírias, jargões, frases feitas, ditados populares, termos regionais ).
Assim, espera -se que o participante respeite a adequação vocabular não usando gírias ou expressões coloquiais, evite repetição desnecessária de palavras e utilize um vocabulário mais formal, como solicitado por um texto dissertativo.
Observações : - A inteligibilidade relacionada ao atendimento das exigências de estruturação textual foi avaliada na segunda competência ( relativa aos aspectos textuais ). - Os problemas de coerência textual provocados por uso indevido do vocabulário foram avaliados na terceira competência ( relativa aos aspectos morfossintáticos e vocabulares ). - Não foram penalizados como desvios morfossintáticos os problemas de caligrafia ( -a/ -o, -s, -r ). - Não foram penalizados os textos grafados integralmente em caixa alta, já que não seria possível distinguir alguma marcação especial para as letras em início de frase. - Cada desvio foi considerado como uma ocorrência, mesmo que dois desvios fossem relativos ao mesmo aspecto linguístico.
A grade de avaliação do desempenho linguístico considerou, portanto, três grandes grupos de competências, segundo os aspectos explicitados anteriormente : a ) domínio das convenções ortográficas : grafia de vogais e consoantes, uso de maiúsculas e minúsculas, emprego do hífen e acentuação gráfica ; b ) domínio dos procedimentos de estruturação textual do ponto de vista microestrutural : organização interna dos períodos, emprego de conectores para a articulação lógica entre os períodos e entre os parágrafos, emprego de marcas de referenciação lexical e pronominal ; utilização dos sinais de pontuação que contribuem para a organização lógica da frase ; c ) domínio das regras de caráter morfossintático estabelecidas como modelares do ponto de vista da modalidade escrita formal da norma-padrão da Língua Portuguesa : concordância nominal e verbal, regência nominal e verbal, colocação pronominal, flexão nominal e verbal, correlação entre tempos e modos verbais, ausência de marcas de oralidade.
A seleção vocabular adequada à modalidade escrita formal da Língua Portuguesa foi incorporada a essa última competência, tendo em vista a intersecção entre as duas do ponto de vista das exigências do registro formal da modalidade escrita da norma-padrão.
Apreciação geral do desempenho dos estudantes : A correção revelou desempenhos distintos dos participantes, marcados pela falta de repertório cultural da maioria dos participantes : - a questão 1 propiciou a oportunidade de maior desenvolvimento do tema solicitado, gerando consequências na estruturação textual, já que os textos foram mais longos, mais elaborados e fluentes.
Entretanto, em virtude da diversidade de abordagens dos textos motivadores ( a falta de penicilina, a recusa no uso da camisinha pelos homens, a maior frequência das mulheres na realização de exames ginecológicos, a transmissão das DST aos bebês durante a gestação ), observaram -se respostas que se configuraram como verdadeiras paráfrases, com comprometimento do caráter autoral desejável em uma questão discursiva de Formação Geral.
Alguns participantes selecionaram um trecho de cada texto motivador para compor suas respostas, outros copiaram trechos inteiros. - a questão 2 teve um comportamento distinto, devido ao caráter polêmico que o caracteriza.
Assim, os textos são mais curtos, com vocabulário repetitivo e preso aos textos motivadores.
Além disso, observou -se grande ocorrência de protesto em função do tema, ora por motivos religiosos ora por preconceito de gênero.
Quanto aos aspectos linguísticos analisados durante esta avaliação, os resultados observados estão descritos nos parágrafos seguintes.
Aspectos ortográficos : O desempenho dos participantes revelou uma diferença muito grande nos dois aspectos analisados nesta competência : baixo índice de desvios da grafia padrão e grande índice de desvios de acentuação.
Em vários casos, ocorre ausência completa de acentuação gráfica.
Os resultados revelam que a tendência dominante entre os universitários brasileiros é a eliminação da acentuação gráfica, talvez motivada pelos hábitos relacionados às redes sociais e pela ausência de esclarecimento dos meios de comunicação, das autoridades e das escolas sobre as decisões do Acordo Ortográfico da Língua Portuguesa de 1990.
Os casos mais sistemáticos de eliminação do acento indicador da sílaba tônica são :  palavras proparoxítonas ( “ sifilis ”, “ proximos ”, “ politicas ”, “ publicas ”, “ transgeneros ” ) ;  palavras paroxítonas terminadas em ditongo crescente ( “ necessario ”, “ noticiarios ”, “ individuo ”, “ dependencia ”, “ varios ”, “ propria ”, “ transmissiveis ”, “ ocorrencia ” ) ;  palavras paroxítonas com hiato ( “ saude ” )  palavras oxítonas ( “ ninguem ”, “ esta ”, “ ate ”, “ tambem ”, “ prevencao ” ).
Por outro lado, destaca -se o uso indevido do acento gráfico em palavras como “ genêro ”, “ melâncolia ”, “ prevênção ”, “ intervênção ”, por exemplo.
Há, também, uma tendência a acentuar paroxítonos como se fossem oxítonos : “ transmissivéis ”.
Quanto ao domínio das convenções relativas à grafia das palavras, observam -se desvios recorrentes, como a hipercorreção pela escolha de “ e ” no lugar de “ i ”, por influência de hábitos da oralidade ( “ descriminação ” no lugar de “ discriminação ”, “ entervenção ” no lugar de “ intervenção ” ) ; desvios de grafia relacionados à variação diastrática podem ser observados em “ estrupo ”, “ subjulgadas ”, “ precoseito ”, “ soubre ”, “ vecendo ”, “ indesencia ”, “ apolojia ”, “ fulga ”, “ dereitos ”, “ sifelis ”.
Observam -se, também, casos de inadequação no uso da maiúscula para destacar determinadas palavras-chave do texto, como “ Transgêneros ”, “ Brasileiros ”, “ Homens ”.
Destaque -se, também, o grande número de participantes que grafam os textos inteiramente em caixa alta.
Vale observar, também, que, ao contrário do que se esperava, não apareceram abreviaturas próprias do “ internetês ”, relacionadas ao uso de redes sociais e e-mails.
Aspectos textuais : Esta competência é a que se revela como a mais problemática entre os participantes, porque são muitos os problemas observados, desvios acumulados durante toda a formação do estudante e que não se resolvem com um estudo autodidata, como acontece com regras ortográficas ou morfossintáticas.
São eles : sequência justaposta de ideias sem encaixamentos sintáticos ; redução drástica de estruturas subordinadas, ao lado do aumento na frequência de estruturas coordenadas e absolutas ; redução no uso de conectores para expressar relações lógicas essenciais à construção do texto, substituídas pela exigência de inferência por parte do interlocutor para suprir a sua ausência ; emprego equivocado de operadores que não estabelecem relações lógicas coerentes entre ideias do texto ; emprego inadequado do pronome relativo ( com omissão da preposição ou a utilização de pronome inadequado, como “ onde ” ) ; repetição exaustiva de termos sem a utilização de procedimentos mais sofisticados de substituição ( hiperonímias, hiponímias, nominalizações, expressões metafóricas ) ; frases fragmentadas que comprometem a estrutura lógico-gramatical ; frases formadas apenas por oração subordinada, sem oração principal.
Um importante aspecto a destacar é o baixo desempenho de uma parte dos participantes em relação à estrutura formal do texto produzido, o que é preocupante ao se levar em conta que são graduandos em fase final de formação.
São frequentes os casos de desvios de estruturação frasal, com uso inadequado ou ausência de conectivos entre parágrafos e entre frases.
Em uma parte dos textos, falta textualidade e domínio do registro padrão da língua.
Na verdade, observam -se relações linguísticas quase agramaticais, como as estabelecidas pela sequência de gerúndios sem o apoio de um ponto de partida para a organização das informações gramaticais e semânticas.
Observou -se que uma parte dos participantes não distribuiu as ideias em parágrafos, talvez devido ao pequeno número de linhas disponibilizadas para a resposta da questão ou, quem sabe, pela suposição de que não seria necessária essa divisão por não se tratar de um texto no modelo de uma redação dissertativo-argumentativa, como solicitado nos vestibulares.
Em função do encaminhamento dos enunciados das duas questões, que solicitaram uma análise do problema e encaminhamentos de políticas públicas, houve uma grande tendência, também, de construção de dois parágrafos desconexos, sem utilização de elementos coesivos adequados à progressão textual.
Quanto à utilização dos mecanismos de referenciação, deve -se destacar a ocorrência de repetições de palavras ou expressões sem a utilização de termos sinônimos ou pronomes, como seria adequado.
Quanto à utilização dos sinais de pontuação, observou -se uma grande precariedade nos textos analisados.
É muito frequente a ocorrência de parágrafos sem marca interna de pontuação para separar os períodos.
Vale observar que não foi penalizada a ausência de vírgula para destacar locuções ou adjuntos adverbiais de pequena extensão deslocados de posição na frase, por ser um uso opcional.
São os seguintes os tipos de problemas encontrados : a ) vírgula : utilização de vírgula para separar o sujeito e o predicado ; ocorrência de apenas uma das vírgulas para separar uma palavra, uma expressão ou uma oração encaixada ; uso de vírgula no lugar do ponto para separar ideias que constituem períodos distintos ; ausência de vírgula para separar elementos de uma enumeração ; ausência de vírgula para separar oração adjetiva explicativa ou utilização inadequada para separar oração adjetiva restritiva ; b ) ponto e vírgula : utilização do ponto e vírgula no lugar de vírgula ; c ) ponto final : ausência de ponto final para separar períodos.
Aspectos morfossintáticos e vocabulares : Em relação à regência, o desvio mais frequente é a falta do sinal indicativo da crase – isso revela que o usuário não tem consciência de que, sob a forma do termo “ a ”, existe a presença de uma contração entre a preposição “ a ” ( exigida pela regência do termo anterior ) e o artigo definido “ a ”.
Um desvio de regência significativo, nos últimos anos, é a utilização inadequada de uma preposição ou sua ausência após o verbo ou o nome ( substantivo ou adjetivo ).
Outro problema relacionado à regência verbal e à nominal, encontrado frequentemente nas questões, foi a ausência de preposição antes de pronome relativo, processo generalizado na modalidade oral da língua, em situações de registro informal.
Apesar da possibilidade de que essa alteração de regência se generalize no padrão escrito da Língua Portuguesa, como já está ocorrendo até em textos jornalísticos, o não emprego da preposição foi penalizado neste processo de avaliação.
A concordância verbal e a concordância nominal apresentam alguns desvios muito frequentes.
Quanto à concordância de número, observou -se ausência de marca ( com sujeito anteposto ou posposto ) ou uso indevido ( uso inadequado da marca de plural comandado pelo núcleo plural da locução adjetiva, apesar de o substantivo que funciona como núcleo do sintagma nominal estar no singular ).
Uma ocorrência que se destacou foi a ausência de acento circunflexo na forma plural do presente do indicativo do verbo “ ter ”, que foi considerada como um desvio na concordância verbal e não na acentuação gráfica.
Quanto à concordância de gênero, vários casos foram observados, normalmente no âmbito de sintagmas nominais longos, em que o adjetivo está afastado do substantivo.
Deve -se destacar o aparecimento da marca de plural em verbos ou adjetivos comandados por adjunto adnominal plural, apesar de serem relacionados a núcleos substantivos no singular, evidenciando um processo de hipercorreção ( exemplo : “ O nome social dos trangêneros podem ser estabelecidos por uma legislação específica ” ).
Quanto à questão da colocação pronominal, foram poucos os desvios observados.
Concluiu -se que, no registro escrito formal, a maioria dos participantes já incorporou regras como a não introdução da frase por um pronome oblíquo e a próclise na presença de um termo atrator.
Não se adotou, entretanto, o padrão excessivamente formal descrito pelas gramáticas normativas em relação à posição do pronome oblíquo em locuções verbais, já que esse uso está muito distante da prática cotidiana, até em textos mais formais.
Quanto aos aspectos vocabulares, alguns tipos de inadequação foram observados : expressões da oralidade ; seleção vocabular incompatível com o contexto, gerando falta de inteligibilidade ; falta de domínio de vocabulário mais abstrato e de maior complexidade, essencial ao desenvolvimento do texto de base dissertativa.
O principal aspecto observado foi a excessiva repetição de certas palavras, revelando limitação de repertório vocabular.
O termo “ pessoa ”, por exemplo, chega a ser repetido até 7 ou 8 vezes em um mesmo texto.
Várias marcas de oralidade foram identificadas, embora não com alta frequência : uso do pronome relativo “ onde ” como relativo universal, falta de artigo definido antes de substantivo, reduções como “ tá ”, “ pra ”, “ pro ”, “ prum ”, expressões informais, eliminação de preposições.
Em função do tema solicitado na questão 1, a sigla DST foi grafada de diferentes maneiras : DST´s, DSTs ou DSTS.
Além disso, houve flutuação de gênero gramatical nos sintagmas que envolviam a sigla : “ os diferentes DSTs ” ou “ as diferentes DSTs ”.
Em função do tema solicitado na questão 2, houve flutuação de gênero gramatical nos sintagmas que envolviam a sua palavra-chave : “ os trans ” ou “ as trans ” ; “ os transgêneros ” ou “ as transgêneras ”.
Todas essas formas foram aceitas, por considerarmos que sua utilização é muito recente e os participantes ainda não sistematizaram essas formas em seu uso da língua. 6.3.2 Componente de Conhecimento Específico Na parte da prova relativa às questões discursivas do Componente de Conhecimento Específico ( Tabela 6.17 ), observa -se que a Média foi bem mais baixa do que para as questões discursivas do Componente de Formação Geral.
Enquanto no Componente de Formação Geral a Média para estudantes de Engenharia de Controle e Automação de todo o Brasil foi 54,5, na parte de Conhecimento Específico a Média foi 11,1.
A maior Média deste componente foi obtida pelos estudantes da região Sul ( 18,3 ), e a menor, pelos da região Norte ( 7,9 ).
Quanto à variabilidade das notas, o Desvio padrão de todo o Brasil foi 16,9.
O maior Desvio padrão foi encontrado na região Sul ( 21,2 ), e o menor, na região Norte ( 14,5 ).
A maior nota Máxima, 96,7, foi obtida nas regiões Sudeste e Sul.
Nas demais regiões a nota Máxima foi : 63,3 na região Norte, 90,0 na região Nordeste e 95,0 na região Centro-Oeste.
A nota Mínima ( 0,0 ) foi obtida por pelo menos um aluno em todas as regiões do Brasil.
A Mediana do Brasil como um todo foi 0,0, a mesma para quatro das cinco regiões, com exceção da região Sul, que obteve Mediana 10,0.
Isso indica que pelo menos metade dos estudantes dessas quatro regiões obteve nota zero nas três discursivas do Componente Específico da área.
O Gráfico 6.10 representa a distribuição das notas nas questões discursivas no Componente de Conhecimento Específico.
A moda desta distribuição ocorre no intervalo [ 0 ; 10 ], com mais do que 65 \% do total de participantes.
Observa -se que nesse intervalos estão computadas as ocorrências de respostas em branco para as três questões discursivas de Conhecimentos Específicos.
A análise de cada uma destas questões será feita a seguir. 6.3.2.1 Análise da Questão Discursiva 3 do Componente de Conhecimento Específico Na questão 3, cujos resultados aferidos encontram -se descritos na Tabela 6.18, a Média dos estudantes de todo o Brasil foi 5,2, a questão de mais baixo desempenho dentre as três discursivas de Conhecimentos Específicos.
A menor Média nessa questão foi obtida pelos alunos da região Nordeste ( 4,0 ), enquanto a maior Média foi obtida na região Sul ( 8,4 ).
Quanto à variabilidade das notas, o Desvio padrão de todo o Brasil foi 18,3.
O maior Desvio padrão foi obtido na região Centro-Oeste ( 24,3 ), enquanto o menor foi obtido na região Nordeste ( 16,6 ).
A nota Máxima, 100,0 pontos, foi alcançada por, pelo menos, um aluno de todas as regiões.
A Mediana do Brasil como um todo foi 0,0, a mesma para todas as regiões, indicando que pelo menos a metade dos estudantes de todas as regiões, sem exceção, obteve grau zero nessa questão.
A nota Mínima ( 0,0 ) também foi a mesma em todas as regiões do Brasil.
O Gráfico 6.11 mostra a distribuição das notas na questão discursiva 3, do Componente de Conhecimento Específico da área de Engenharia de Controle e Automação.
A moda dessa distribuição foi a resposta em branco ( mais de 50 \% dos participantes ) seguida da nota zero para cerca de 36 \% dos estudantes, compatível com o resultado de Mediana zero em todas as regiões. 6.3.2.2 Comentários sobre as respostas à Questão Discursiva 3 A Questão 3 abordava um conteúdo básico na área de Controle e Automação : determinar o erro em regime permanente ( estacionário ) para um sistema de controle em malha fechada ( controle de um motor CC ).
Este é um problema usual na literatura empregada nas disciplinas relacionadas ao tema.
Sendo um conteúdo básico da área de controle, a questão poderia ter sido classificada como fácil ou, no máximo, como média.
A correção evidenciou que alguns estudantes confundiram controle de posição com constante de posição ‘ Kp ’ ( para uma entrada degrau unitário ).
Isso pode ser efeito do próprio enunciado, que menciona"o controle de posição do motor", sendo que a entrada é do tipo rampa unitária e a constante de ‘ Kv ’ é nomeada constante de velocidade.
Para alguns respondentes, isso pode ter causado dúvidas para a solução da questão.
A maioria das respostas revelou baixo ou completo desconhecimento a respeito do problema proposto e dos termos técnicos específicos de controle.
No entanto, apesar de ser uma minoria, alguns demonstraram pleno conhecimento e responderam perfeitamente à questão.
Entre as resposta encontradas, muitas tentavam utilizar a definição básica do erro estacionário, entrada menos a saída, quando ‘ t ’ tende ao infinito ou ‘ s ’ tende a zero ( lim s E(s) ).
Vários respondentes iniciaram com erro no cálculo da função de transferência de malha fechada e, por conseguinte, erraram o resto da questão.
Dos mais de 4.800 formandos que participaram do exame, em torno de 2.300 tentaram resolver a questão, e aproximadamente 2.500 deixaram a questão em branco.
Dentre os que responderam, a maioria absoluta ( em torno de 2.100 ) teve notas menores do que 3,0, mostrando claramente que os respondentes têm pouco ou nenhum conhecimento do assunto.
Uma minoria, em torno de 230 respondentes, mostrou algum conhecimento ( notas acima de 4,0 ), dos quais pouco mais de 130 respondentes acertaram totalmente a questão e obtiveram nota 10,0.
Sendo uma questão básica de Controle e Automação, é preocupante verificar que o tipo de formação com a qual estão saindo das universidades, não permita, minimamente, resolver um problema elementar que, certamente, enfrentarão no dia a dia do exercício profissional.
Foi observado que, para os respondentes em geral, a questão foi difícil, pois ou erraram no desenvolvimento ou não tinham qualquer conhecimento do conteúdo. 6.3.2.3 Análise da Questão Discursiva 4 do Componente de Conhecimento Específico A Tabela 6.19 contém as informações relativas à questão 4 do conjunto de questões do Componente de Conhecimento Específico.
O desempenho dos estudantes nessa questão foi o melhor dentre as três questões discursivas.
A Média geral do Brasil foi 19,4, sendo a menor Média registrada na região Norte ( 12,8 ), e a maior, na região Sul ( 30,6 ).
A nota Máxima ( 100,0 ) foi atingida por, pelo menos, um concluinte de todas as regiões.
A Mediana em todo o Brasil foi zero, com exceção da região Sul ( 10,0 ), indicando que, com exceção dessa região, nas demais, pelo menos 50 \% dos participantes obteve nota zero nessa questão.
A nota Mínima de todas as regiões foi zero.
O Gráfico 6.12 representa a distribuição de notas da questão discursiva 4, do Componente de Conhecimento Específico.
Essa distribuição tem moda na classe de estudantes que deixou a resposta em branco ( cerca de 54 \% ). 6.3.2.4 Comentários sobre as respostas à Questão Discursiva 4 O enunciado da Questão 4 exigia conhecimentos básicos e fundamentais a respeito de Sistemas de Controle de Aquecimento.
O enunciado era claro, com informações e dados mais do que suficientes para a resposta.
Trata -se de uma questão de solução rápida, necessitando de pouco espaço para desenvolvimento da resposta, pressupondo a objetividade.
A questão pode ser considerada fácil, porque trata de conhecimentos básicos e fundamentais de instrumentação.
Pelo que foi proposto como padrão de resposta, exigia -se objetividade, precisão e clareza de terminologia.
Não foram verificadas interpretações diferentes do que o padrão esperava e também não foi verificado nenhum caso em que o espaço para resposta fosse insuficiente, mesmo quando a resposta perdeu objetividade.
Foram observados os seguintes casos : respostas que seguiram o padrão integralmente ; as que seguiram o padrão, mas não completaram a resposta ou continham algum erro ; as que parafrasearam termos do enunciado sem apresentar consistência na resposta. 6.3.2.5 Análise da Questão Discursiva 5 do Componente de Conhecimento Específico A Tabela 6.20 contém as informações relativas à questão 5 do conjunto do Componente de Conhecimento Específico.
O desempenho dos estudantes de todo o Brasil nesta questão foi mais elevado que o da questão 3, mas muito abaixo ao da questão discursiva 4.
A nota Média dos estudantes de todo o Brasil foi 8,8.
A maior Média foi registrada resposta.
Houve, também, respostas em que foram apenas"jogados"termos e números de forma incoerente.
De modo geral, a linguagem adotada pelos formandos foi adequada aos termos específicos da Área, considerando -se as que continham acertos parciais ou totais.
A abordagem mais frequente para responder à questão foi seguir a ordem das perguntas e manter a objetividade e precisão que o enunciado sugeria.
Ficaram evidentes erros relacionados à falta de conhecimento de conceitos - que causaram, em muitos casos, respostas equivocadas - ; de terminologia de grandezas e medidas ; e de avaliação \textbf{numérica} de precisão de leitura de sensores.
Foram também observadas respostas que pareciam aleatórias, mas que talvez tivessem apoio em alguma noção imprecisa do tema da questão.
Como na Questão 3, a quantidade de respostas deixadas em branco foi muito alta.
De aproximadamente de 4.800 participantes, aproximadamente 2.200 tentaram fazer a questão, e quase 2.600 deixaram a questão em branco.
Em relação às questões respondidas, quase 1.000 obtiveram notas menores do que 3,0, mostrando que muitos formandos têm pouco ou nenhum conhecimento sobre o assunto.
O restante ( em torno de 1.200 respondentes ) mostrou algum conhecimento sobre o tema abordado : metade com notas entre 4,0 e 6,0 e a outra metade com notas maiores do que 7,0.
A quantidade de respostas que obteve nota máxima não passou de 130.
Novamente, o desempenho dos estudantes pode ser considerado insuficiente, tendo em vista o grau de dificuldade da questão, a quantidade de respostas incompletas e os tipos de erros cometidos.
Dado a pouca dificuldade da questão e o fato de abordar fundamentos, pode -se considerar que existem deficiências de formação em relação ao conteúdo exigido.
As mais marcantes são : falta de solidez no uso de nomenclatura e conceitos de controle de processos e em conhecimentos relativos a grandezas de medidas, propriedades físicas e noções metrológicas associadas a sensores ( faixa de medição, resolução, precisão, etc. ) analógicos e digitais.
São conceitos que, se não forem suficientemente assimilados e consolidados, podem comprometer o exercício da profissão. 6.3.2.5 Análise da Questão Discursiva 5 do Componente de Conhecimento Específico A Tabela 6.20 contém as informações relativas à questão 5 do conjunto do Componente de Conhecimento Específico.
O desempenho dos estudantes de todo o Brasil nesta questão foi mais elevado que o da questão 3, mas muito abaixo ao da questão discursiva 4.
A nota Média dos estudantes de todo o Brasil foi 8,8.
A maior Média foi registrada na região Sul ( 15,7 ), enquanto a menor Média foi registrada na região Norte ( 4,8 ).
Quanto à variabilidade das notas, o Desvio padrão dos alunos do Brasil, como um todo, foi 17,3.
Enquanto o maior desvio foi encontrado na região Sul ( 23,6 ), o menor foi encontrado na região Norte ( 13,0 ).
Para o conjunto de alunos de Engenharia de Controle e Automação do Brasil, a nota Máxima foi 100,0 em todas as regiões, com exceção da região Norte ( 95,0 ), e a nota Mínima foi zero em todas as regiões.
A Mediana também foi zero para o Brasil e em todas as regiões, indicando que pelo menos a metade dos alunos recebeu esta nota, em cada região.
O Gráfico 6.13 representa a distribuição das notas da questão discursiva 5 do Componente de Conhecimento Específico.
Essa distribuição, como todas as relativas às questões discursivas dessa parte da prova, também tem moda principal nas respostas em branco.
Destaca -se que a frequência de notas zero somada à de estudantes que não responderam à questão chega a, aproximadamente, 70 \% do total de participantes. 6.3.2.6 Comentários sobre as respostas à Questão Discursiva 5 O enunciado da Questão 5 era claro e objetivo.
O conteúdo, a abordagem e a linguagem empregada foram adequados e integralmente aderentes às Diretrizes da Engenharia de Controle e Automação.
O nível de profundidade exigido é básico na área de Robótica.
O padrão de resposta foi apresentado em linguagem adequada e de acordo com a terminologia e procedimentos das respectivas áreas.
A questão era bastante objetiva, solicitava um cálculo no item ‘ a ’, a partir dos dados disponíveis e, no item ‘ b ’, a especificação das equações da cinemática direta para o manipulador que descrevem o movimento do plano horizontal em função dos ângulos das duas primeiras articulações ( 1 e 2 ).
Assim, no item ‘ a ’, a menos de erros de cálculos, não se esperava resposta diferente da prevista no padrão.
O item ‘ b ’, que era mais difícil, envolvia conhecimento do modelo físico do robô SCARA e algumas manipulações algébricas.
Para este item o padrão apresentava dois caminhos de resposta, uma a partir das relações geométricas do mecanismo que representa o robô e outra a partir dos \textbf{parâmetros} de Denavit-Harteenberg. 6.3.3 Considerações Finais De maneira geral, fica evidente que os futuros engenheiros não estão conseguindo apreender o ferramental fundamental da área como se esperaria.
Os resultados apontam no sentido de que os estudantes não estão sendo expostos de forma adequada/correta aos Não foram encontradas respostas corretas diferentes das apresentadas pelo padrão proposto.
Aqueles que responderam integralmente ou mesmo parcialmente, o fizeram de acordo com o encaminhamento que o padrão apresentava.
Para o item ‘ a ’ as respostas corretas e as parcialmente corretas foram dadas identificando corretamente a sequência de operações e tempos associados, para finalmente especificar a quantidade de robôs necessária.
Já com relação ao item ‘ b ’, poucos chegaram até às equações solicitadas, alguns identificaram o modelo físico do robô SCARA e desenvolveram parcialmente o modelo cinemático, a maioria através das relações geométricas e poucos empregando as matrizes de transformação baseadas nos \textbf{parâmetros} de Denavit-Hartenberg.
Para os que tentaram responder o item ‘ b ’ de forma errada, ficou claro que o fizeram por desconhecimento do procedimento de modelagem de manipuladores robóticos.
Novamente, num universo de pouco mais de 4.800 respondentes, 2.200 tentaram responder à questão, os demais a deixaram em branco.
Dentre os que responderam, em torno de 1.600 obtiveram notas menores do que 3,5.
Cerca de 10 \% dos participantes demostrou algum conhecimento sobre os conteúdos envolvidos na questão, sendo que uma minoria acertou totalmente a questão e obtendo nota 10,0.
Portanto, para os participantes, em geral, a questão foi difícil.
Apesar de não ser uma questão básica de Controle e Automação, verificou -se que os estudantes de Engenharia de Controle e Automação estão concluindo o ensino superior sem os conhecimentos de Robótica avaliados pela Questão 5, conhecimentos fundamentais para um profissional da Área, o que é, no mínimo, preocupante.
Os formandos não estão sendo expostos ao conteúdo de Robótica tratado na questão, nem nos seus aspectos práticos ( item ‘ a ’ ) e, muito menos, nos mais teóricos ( item ‘ b ’ ).
Apesar de não ser um conteúdo básico, é fundamental para o exercício profissional, principalmente atualmente, considerando o aumento da robotização dos sistemas em geral.
Saber analisar o desempenho de um manipulador robótico, e desenvolver seu modelo matemático, visando as simulações e o projeto, são quesitos necessários para os engenheiros de Controle e Automação, e, pelo visto, não está sendo dada a devida importância a esta área nos cursos de formação. 6.3.3 Considerações Finais De maneira geral, fica evidente que os futuros engenheiros não estão conseguindo apreender o ferramental fundamental da área como se esperaria.
Os resultados apontam no sentido de que os estudantes não estão sendo expostos de forma adequada/correta aos conteúdos elementares, nem nos seus aspectos teóricos ( Questão 3 e Questão 5, item ‘ b ’ ), mas também nos mais práticos ( Questão 4 e Questão 5, item ‘ a ’ ).
Esse alerta, dado pelo Enade, é indicativo da necessidade urgente de se rever a maneira pela qual tais conteúdos são transmitidos e como é avaliada o seu aprendizado pelos respondentes.
Caso contrário, teremos no Brasil profissionais cada vez menos qualificados e mais despreparados para os reais desafios da engenharia, sob pena de perdermos, definitivamente, capacidade de desenvolvimento tecnológico de alto nível e internacionalmente competitivo.

% matched lemmas: numérico, parâmetro
\end{textsample}
